\documentclass[10pt,a4paper]{amsart}
\usepackage[margin=19mm]{geometry}
\usepackage{amsmath,amssymb,mathtools}
\usepackage[scr=rsfs,cal=euler]{mathalfa}
\usepackage{xfrac}
\usepackage[unicode,hidelinks,bookmarks=false]{hyperref}
\newcommand{\ie}{i.e.\ }
\newcommand{\cf}{cf.\ }
\newcommand{\resp}{respectively\ }
\newcommand{\Subsection}{\subsection}
\providecommand{\nobreakdash}{\nobreak\hbox{-}}
\setlength{\emergencystretch}{2em}
\multlinegap0pt
\usepackage{bm}
\usepackage{bbm}
\usepackage{dsfont}
\usepackage{xspace}
\usepackage{cancel}
\usepackage{mathtools}
\usepackage{amsthm}
\makeatletter
\@ifundefined{theorem}{\newtheorem{theorem}{Theorem}}{}
\@ifundefined{proposition}{\newtheorem{proposition}[theorem]{Proposition}}{}
\makeatother
\newcommand{\Q}{\mathbb{Q}}
\newcommand{\Z}{\mathbb{Z}}
\title[Counting rational maps on $\mathbb{P}^1$ with prescribed loc]{Counting rational maps on $\mathbb{P}^1$ with prescribed local conditions}
\author{Khoa D. Nguyen, Anwesh Ray}
\date{Source: arXiv:2408.15648v3 (CC BY 4.0)}
\begin{document}
\maketitle

\noindent\emph{Source and licence.} Counting rational maps on $\mathbb{P}^1$ with prescribed local conditions. Version arXiv:2408.15648v3 (CC BY 4.0), distributed under the Creative Commons Attribution 4.0 International licence, as recorded in the arXiv licence metadata for this exact version and shown on the abstract page https://arxiv.org/abs/2408.15648v3. Full rights notice: licenses/arxiv-2408.15648-CC-BY-4.0.txt.\par\smallskip

\noindent\textit{Editorial scope.} Counted unit: the Proposition whose source label is \texttt{prop:barBp'} in the article (source0.tex lines 1152--1154, source label \texttt{prop:barBp'}), delivered together with the article's own complete original proof of it (source0.tex lines 1155--1187, one \texttt{proof} environment of 33 lines). No mathematics is written, repaired or re-derived: statement and proof are the article's own text line for line. Required context retained uncounted: none. Every definition, display and label of those lines is retained unchanged. Omitted from this package and not redistributed: 1--1151, 1188--1201. Nothing inside a retained excerpt is omitted or altered: every retained body is delivered exactly as the article's lines are written, so no omission receipt of any kind accompanies this package. Citations: the retained excerpts carry no citation command at all, so no bibliography entry of the article is transcribed and no reference is redistributed. Reference adaptation: none was needed and none was made; every reference inside the delivered unit resolves inside it, so no unresolved reference or citation is produced. Printed slips kept literal: no printed word, symbol, hypothesis or number of the counted statement or of its proof was altered. Layout and class: the article's own class is \texttt{amsart} with packages including geometry, amscd, amsmath, amsxtra, amsthm, amssymb, stmaryrd, xr, mathrsfs, mathtools, enumerate, commath, comment, multirow; the wrapper is the lane's pinned \texttt{amsart} editorial wrapper at 10pt on A4, which restates the theorem-like environments the retained text needs and adds the article's own macro definitions that the retained text needs (\texttt{\textbackslash Q}, \texttt{\textbackslash Z}), with extra wrapper packages (bm, bbm, dsfont, xspace, cancel, mathtools). The delivered numbering is the unit's own. Assets: no retained span contains a figure environment, a graphics-inclusion command, a PDF-page inclusion, a table or a TikZ picture; no image file is copied into this package, the package declares no asset dependency and no image file is redistributed with it. Licence: version arXiv:2408.15648v3 of this article is distributed under the Creative Commons Attribution 4.0 International licence, as recorded in the arXiv licence metadata for this exact version and shown on the abstract page https://arxiv.org/abs/2408.15648v3. A full rights notice is delivered at licenses/arxiv-2408.15648-CC-BY-4.0.txt. Characters: every retained excerpt is pure ASCII, so the delivered engine, XeLaTeX with its default Latin Modern text font, reports no missing character.
\begin{proposition}\label{prop:barBp'}
    When $p>3$, we have $|\bar{B}_p'|<p^4+3p^3$.
\end{proposition}

\begin{proof}
    Let $I$ be the ideal of $\Z[a,\ldots,f]$ generated by $R$ and its partial derivatives. By using Macaulay2\footnote{Available at \url{https://macaulay2.com/}}, we can compute the following primary decomposition in $\Q[a,\ldots,f]$:
    $$I\Q[a,\ldots,f]=(ce-bf,cd-af,bd-ae)\cap (e^2-4df,2cd-be+2af,b^2-4ac).$$

    Let $I_1=(ce-bf,cd-af,bd-ae)$ and $I_2=(e^2-4df,2cd-be+2af,b^2-4ac)$ as ideals of $\Z[a,\ldots,f]$. Using Macaulay2 again, we can check that:
    $$I_1\cap I_2\subseteq I\ \text{and}\ 12I\subseteq I_1\cap I_2.$$
    This means that inside the ring $Z[\frac{1}{12},a,\ldots,f]$, we have $I=I_1\cap I_2$. Since $p>3$, we have $\bar{B}_p'=\bar{C}_p\cup\bar{D}_p$ where
    $$\bar{C}_p:=\{(a,b,c,d,e,f)\in (\Z/p\Z)^6:\ ce-bf=cd-af=bd-ae=0\}\ \text{and}$$
    $$\bar{D}_p:=\{(a,b,c,d,e,f)\in (\Z/p\Z)^6:\ e^2-4df=2dc-be+2af=b^2-4ac=0\}.$$

    It is easy to determine $|\bar{C}_p|$: we simply determine the number of ordered pairs of linearly dependent vectors $((a,b,c),(d,e,f))$ in $(\Z/p\Z)^3$. When $(a,b,c)$ is zero, $(d,e,f)$ can be any vector. Otherwise, for any given non-zero $(a,b,c)$, there are exactly $p$ many choices for $(d,e,f)$. This gives
    \begin{equation}\label{eq:barCp}
    |\bar{C}_p|= p^3+p(p^3-1).
    \end{equation}
    We now give an upper bound for $|\bar{D}_p\setminus\bar{C}_p|$. 

    \textbf{Claim:} for $(a,b,c,d,e,f)\in \bar{D}_p\setminus\bar{C}_p$, we have $abcdef\neq 0$. We present the proof that $a\neq 0$, then the  proof for each of the statements $c\neq 0$, $d\neq 0$, and $f\neq 0$ is
    similar. After proving $acdf\neq 0$, since $e^2=4df$ and $b^2=4ac$ we have $be\neq 0$ as well.
    Suppose $a=0$, then from the definition of $\bar{D}_p$, we have:
    $$b=0,\ dc=0,\ e^2=4df$$
    Since $(a,b)=(0,0)$ we have $c\neq 0$, otherwise $(a,b,c,d,e,f)\in\bar{C}_p$. Hence $d=0$ and $e=0$. But then we still have $(a,b,c,d,e,f)\in\bar{C}_p$, contradiction.

    Now there are $(p-1)^3$ many tuples $(a,b,e)$ with $abe\neq 0$. We fix any such $(a,b,e)$. Then
    $c$ is determined uniquely from $b^2=4ac$. We now have the system of equations in $(d,f)$:
    $$df=e^2/4\ \text{and}\ 2dc-be+2af=0.$$
    By substituting $f=(be-2dc)/(2a)$ into the first equation, we get a quadratic equation in $d$. This explains why there are at most $2$ solutions $(d,f)$ of the above system.  Therefore
    \begin{equation}\label{eq:barDp}
    |\bar{D}_p\setminus\bar{C}_p|\leq 2(p-1)^3.
    \end{equation}

    From \eqref{eq:barCp} and \eqref{eq:barDp}, we have:
    $$|\bar{B}_p'|=|\bar{C}_p|+|\bar{D}_p\setminus\bar{C}_p|<p^4+3p^3.$$
    \end{proof}

\end{document}
